\chapter{The Laplace Equation}

\section{Steady states}

The Laplace equation
\[ u_{xx} + u_{yy} = 0 \]
describes the steady temperature in a flat plate, the electric potential in a region without charge, and the flow of an ideal fluid. A solution is called a \textbf{harmonic function}. There is no time variable, so we only have boundary conditions. A problem that gives $u$ on the whole boundary is a \textbf{Dirichlet problem}.

\begin{defbox}{Two useful facts (stated)}
\begin{itemize}
  \item \textbf{Maximum principle:} a harmonic function takes its largest and smallest values on the boundary, never inside.
  \item \textbf{Uniqueness:} the Dirichlet problem has only one solution. So if you find a solution by any method, it is \emph{the} solution. (For example, if all four sides of a plate are at $100^\circ$, then $u=100$ everywhere.)
\end{itemize}
\end{defbox}

\section{Rectangle}

Take the rectangle $0<x<a$, $0<y<b$ with $u=0$ on three sides and $u=f(x)$ on the top:
\[ u(0,y) = 0,\quad u(a,y)=0,\quad u(x,0)=0,\quad u(x,b)=f(x). \]

\begin{methodbox}{Laplace equation on a rectangle}
\begin{enumerate}
  \item Put $u = X(x)Y(y)$. Then $\dfrac{X''}{X} = -\dfrac{Y''}{Y} = -\lambda$, so $X'' + \lambda X = 0$ and $Y'' - \lambda Y = 0$.
  \item Zero BCs at $x=0$ and $x=a$ give $\lambda_n = (n\pi/a)^2$ and $X_n = \sin\dfrac{n\pi x}{a}$.
  \item The $Y$ equation is $Y'' = (n\pi/a)^2Y$. Its solutions are $\cosh$ and $\sinh$. Since $Y(0)=0$, we keep $Y_n = \sinh\dfrac{n\pi y}{a}$.
  \item Superpose: $u = \displaystyle\sum_{n=1}^\infty B_n\sinh\frac{n\pi y}{a}\,\sin\frac{n\pi x}{a}$.
  \item At $y=b$: $B_n\sinh\dfrac{n\pi b}{a} = \dfrac2a\displaystyle\int_0^af(x)\sin\frac{n\pi x}{a}\,\dd x$.
\end{enumerate}
\end{methodbox}

Notice the change from Chapters 10--11. The condition on $y$ is a \emph{boundary} value, not an initial value, and we get $\sinh$ and not an exponential decay or a cosine. If the other sides also have non-zero values, solve one problem for each side (rotate the roles of $x$ and $y$) and add the solutions by superposition.

\begin{examplebox}{Square plate, one hot side}
$a=b=\pi$, three sides at $0^\circ$, the top side at $100^\circ$. Then $B_n\sinh n\pi = \dfrac2\pi\displaystyle\int_0^\pi100\sin nx\,\dd x$, which is $\dfrac{400}{n\pi}$ for odd $n$. So
\[ u(x,y) = \frac{400}{\pi}\sum_{n\ \text{odd}}\frac{\sinh ny\,\sin nx}{n\,\sinh n\pi}. \]
\textbf{Centre check.} At $(\pi/2,\pi/2)$: $n=1$ gives $\frac{400}{\pi}\cdot\frac{\sinh(\pi/2)}{\sinh\pi} = 127.32\times0.19926 = 25.37$; $n=3$ gives $-0.381$; $n=5$ gives $+0.020$. The sum is $25.0$. By symmetry the exact answer is $100/4=25$, because four rotated copies of this problem add to $u=100$ everywhere (uniqueness). \checkmark
\end{examplebox}

\section{Disk}

For a disk we use polar coordinates $(r,\theta)$:
\[ u_{rr} + \frac1ru_r + \frac1{r^2}u_{\theta\theta} = 0,\ \ 0\le r<a,\qquad u(a,\theta) = f(\theta). \]

\begin{derivbox}{The Laplace equation in polar coordinates}
Let $x = r\cos\theta$, $y = r\sin\theta$. Then $r = \sqrt{x^2+y^2}$, so $r_x = \cos\theta$, $r_y = \sin\theta$, and $\theta = \arctan(y/x)$, so $\theta_x = -\dfrac{\sin\theta}{r}$, $\theta_y = \dfrac{\cos\theta}{r}$. By the chain rule the two derivative operators are
\[ \partial_x = \cos\theta\,\partial_r - \frac{\sin\theta}{r}\,\partial_\theta,\qquad \partial_y = \sin\theta\,\partial_r + \frac{\cos\theta}{r}\,\partial_\theta. \]
Apply $\partial_x$ twice (remember that $\cos\theta$ and $\sin\theta$ also depend on $\theta$):
\[ \begin{aligned}u_{xx} = {}& \cos^2\theta\,u_{rr} - \frac{2\sin\theta\cos\theta}{r}\,u_{r\theta} + \frac{\sin^2\theta}{r}\,u_r \\ & + \frac{2\sin\theta\cos\theta}{r^2}\,u_\theta + \frac{\sin^2\theta}{r^2}\,u_{\theta\theta}.\end{aligned} \]
The same steps for $\partial_y$ give
\[ \begin{aligned}u_{yy} = {}& \sin^2\theta\,u_{rr} + \frac{2\sin\theta\cos\theta}{r}\,u_{r\theta} + \frac{\cos^2\theta}{r}\,u_r \\ & - \frac{2\sin\theta\cos\theta}{r^2}\,u_\theta + \frac{\cos^2\theta}{r^2}\,u_{\theta\theta}.\end{aligned} \]
Add them. The mixed terms cancel and $\sin^2+\cos^2 = 1$:
\[ u_{xx} + u_{yy} = u_{rr} + \frac1ru_r + \frac1{r^2}u_{\theta\theta}. \]
\end{derivbox}

\begin{derivbox}{The separated equations for the disk}
Put $u = R(r)\Theta(\theta)$ into $u_{rr} + \frac1ru_r + \frac1{r^2}u_{\theta\theta} = 0$ and multiply by $r^2/(R\Theta)$:
\[ \frac{r^2R'' + rR'}{R} = -\frac{\Theta''}{\Theta} = \lambda. \]
So $\Theta'' + \lambda\Theta = 0$ and $r^2R'' + rR' - \lambda R = 0$. Because $\theta$ and $\theta + 2\pi$ are the same point, $\Theta$ must have period $2\pi$. The only periodic solutions come from $\lambda = n^2$ with $n = 0,1,2,\dots$ (for $\lambda<0$ we would get $\cosh$ and $\sinh$, which are not periodic). For the $r$ equation use $R = r^m$: $m(m-1) + m - n^2 = m^2 - n^2 = 0$, so $m = \pm n$. The function $r^{-n}$ is infinite at $r = 0$, so only $r^n$ is allowed for $n\ge1$. For $n=0$ the equation is $r^2R''+rR'=0$, whose solutions are $1$ and $\ln r$, and again only the constant is bounded.
\end{derivbox}

\begin{methodbox}{Laplace equation on a disk}
\begin{enumerate}
  \item Put $u = R(r)\Theta(\theta)$. This gives $\Theta'' + \lambda\Theta=0$ and $r^2R'' + rR' - \lambda R = 0$.
  \item $\Theta$ must be periodic with period $2\pi$, so $\lambda = n^2$ and $\Theta = \cos n\theta,\ \sin n\theta$ ($n = 0,1,2,\dots$).
  \item The $R$ equation is an \textbf{Euler--Cauchy equation} (Chapter 2). With $R=r^m$: $m^2 = n^2$, so $R = r^{n}$ or $r^{-n}$. The value $r^{-n}$ is infinite at the centre, so we keep $r^n$ (and for $n=0$ the constant).
  \item Superpose and match $f(\theta)$ with its Fourier series:
  \[ u(r,\theta) = \frac{a_0}{2} + \sum_{n=1}^\infty\Big(\frac ra\Big)^n\big(a_n\cos n\theta + b_n\sin n\theta\big), \]
  where $a_n$, $b_n$ are the Fourier coefficients of $f(\theta)$ on $(-\pi,\pi)$.
\end{enumerate}
\end{methodbox}

The constant term gives the \textbf{mean value property}: the value at the centre is the average of $f$ around the boundary.

\begin{examplebox}{Disk problems}
\begin{itemize}
  \item $a = 1$, $f(\theta) = \cos^2\theta = \tfrac12 + \tfrac12\cos 2\theta$. Then $u = \tfrac12 + \tfrac12r^2\cos 2\theta = \tfrac12 + \tfrac12(x^2-y^2)$, because $r^2\cos2\theta = x^2-y^2$. It is harmonic: $u_{xx}+u_{yy} = 1 - 1 = 0$.
  \item $a=2$, $f(\theta) = 3 + 4\cos\theta$. Then $u = 3 + 4\dfrac r2\cos\theta = 3 + 2x$.
\end{itemize}
\end{examplebox}

\begin{practicebox}
\begin{enumerate}
  \item Rectangle $0<x<\pi$, $0<y<1$: $u=0$ on the left, right and bottom, $u(x,1)=\sin 2x$. Find $u$.
  \item A square plate has two adjacent sides at $100^\circ$ and the other two at $0^\circ$. What is the temperature at the centre?
  \item Disk of radius 2 with $u(2,\theta) = 3 + 4\cos\theta$: find $u(r,\theta)$ and $u$ at the centre.
  \item Disk of radius 1 with $u(1,\theta) = \sin^2\theta$: find $u(r,\theta)$.
  \item Why is $u = x^2+y^2$ not a solution of the Laplace equation?
\end{enumerate}
\end{practicebox}

\begin{answerbox}
(1) Only $n=2$ is needed: $u = \dfrac{\sinh 2y}{\sinh 2}\sin 2x$. \quad (2) Each hot side contributes $25^\circ$ at the centre by symmetry, so $50^\circ$C. \quad (3) $u = 3 + 2r\cos\theta = 3 + 2x$; at the centre $u = 3$. \quad (4) $\sin^2\theta = \frac12 - \frac12\cos 2\theta$, so $u = \frac12 - \frac12r^2\cos2\theta = \frac12 - \frac12(x^2-y^2)$. \quad (5) $u_{xx}+u_{yy} = 2+2 = 4\neq0$.
\end{answerbox}
